\documentclass[11pt,a4paper]{article}
% Shared preamble for RuPhO English problem statements (match T1 2026 style).
% Use: \input{latex/rupho_en_preamble.tex} after \documentclass.
\usepackage[margin=1in]{geometry}
\usepackage{amsmath,amssymb}
\usepackage{graphicx}
\usepackage{float}
\usepackage{enumitem}
\usepackage{microtype}
\usepackage{caption}
\usepackage{tabularx}
\usepackage{hyperref}

\captionsetup{font=small,labelfont=bf,skip=6pt}

\setlist[itemize]{leftmargin=1.2cm}
\setlist[enumerate]{leftmargin=1.2cm}

% One outer box per subproblem: [id | statement | points] (no inner rules)
\newcommand{\problembox}[3]{%
  \par\addvspace{0.42em}%
  \noindent
  \setlength{\fboxsep}{7.5pt}%
  \setlength{\fboxrule}{0.95pt}%
  \fbox{%
    \begin{minipage}{\dimexpr\linewidth-2\fboxsep-2\fboxrule\relax}
    \begin{tabularx}{\linewidth}{@{}>{\raggedright\arraybackslash\bfseries}p{2.1em} @{\hspace{0.9em}} >{\raggedright\arraybackslash}X @{\hspace{0.9em}} >{\raggedleft\arraybackslash\bfseries}p{2.4em}@{}}
    #1 & #2 & #3
    \end{tabularx}
    \end{minipage}%
  }%
  \par\addvspace{0.32em}%
}


\title{\textbf{Nonideal gas model}}
\author{\normalsize Y14 (2014) --- English translation}
\date{}
\begin{document}
\maketitle

The interaction between the molecules of a substance (liquid) can be described using the dependence of potential energy on the distance between the centers of the molecules. In one of the simplified models, this dependence has the form shown in the figure. The average concentration of molecules is $n$, the mass of a molecule is $m$. The substance consists of monatomic molecules.

\begin{figure}[H]
    \centering
    \includegraphics[width=0.72\textwidth]{T11_EN_assets/fig1.jpg}
    \caption{Figure 1.}
\end{figure}

\problembox{A1}{Determine the internal energy $U$ of one mole of a substance as a function of temperature $T$ and volume $V$. Consider Avogadro's number $N_A$ and Boltzmann's constant $k_B$ given.}{0}

\problembox{A2}{Estimate the molar heat of vaporization (neglect the presence of saturated steam).}{0}

\problembox{A3}{In a vessel of volume $V$ there is $1~\text{mol}$ substances in the gaseous phase under conditions far from critical. As the temperature drops, the condensation process begins. Estimate the temperature $T_к$ at which condensation begins.}{0}

\problembox{A4}{For the case of the liquid phase under conditions far from critical, estimate the surface tension $\sigma$.}{0}

\problembox{A5}{Evaluate the conditions under which a substance can be considered an ideal gas.}{0}

\problembox{A6}{Draw a conclusion about the equation of state of matter (i.e., the $P$, $V$ and $T$ bonds) in the gaseous phase under normal conditions. Note. Use the thermodynamic relation \textbackslash\{\}[\textbackslash\{\}left(\textbackslash\{\}frac\{\textbackslash\{\}partial U\}\{\textbackslash\{\}partial V\}\textbackslash\{\}right)_T = T \textbackslash\{\}left( \textbackslash\{\}frac\{\textbackslash\{\}partial P\}\{\textbackslash\{\}partial T\} \textbackslash\{\}right)_V - P.\textbackslash\{\}]}{0}

Note. Use the thermodynamic relation

\textbackslash\{\}[\textbackslash\{\}left(\textbackslash\{\}frac\{\textbackslash\{\}partial U\}\{\textbackslash\{\}partial V\}\textbackslash\{\}right)_T = T \textbackslash\{\}left( \textbackslash\{\}frac\{\textbackslash\{\}partial P\}\{\textbackslash\{\}partial T\} \textbackslash\{\}right)_V - P.\textbackslash\{\}]

\problembox{A7}{Estimate the mean free path and free travel time of the molecule.}{0}

\vspace{1em}
\noindent\textit{Source:} \url{https://pho.rs/p/3991}

\end{document}